{{دیوان
| عنوان = کُفر و اسلام
| ماخذ = ویکی ماخذ
| ماخذ_لنک = https://ur.wikisource.org/wiki/%DA%A9%D9%8F%D9%81%D8%B1_%D9%88_%D8%A7%D8%B3%D9%84%D8%A7%D9%85
| مقابلہ = اقبال سائبر لائبریری (اقبال اکادمی پاکستان)
| مقابلہ_لنک = https://www.iqbalcyberlibrary.net
}}
=== تضمین بر شعرِ میر رضی دانشؔ ===

<poem>
ایک دن اقبالؔ نے پُوچھا کلیمِ طُور سے
اے کہ تیرے نقشِ پا سے وادیِ سِینا چمن

آتشِ نمرود ہے اب تک جہاں میں شُعلہ ریز
ہوگیا آنکھوں سے پنہاں کیوں ترا سوزِ کُہن

تھا جوابِ صاحبِ سینا کہ مسلم ہے اگر
چھوڑ کر غائب کو تُو حاضر کا شیدائی نہ بن

ذوقِ حاضر ہے تو پھر لازم ہے ایمانِ خلیلؑ
ورنہ خاکستر ہے تیری زندگی کا پیرہن

ہے اگر دیوانۂ غائب تو کچھ پروا نہ کر
منتظر رہ وادیِ فاراں میں ہو کر خیمہ زن

عارضی ہے شانِ حاضر، سطوَتِ غائب مدام
اس صداقت کو محبّت سے ہے ربطِ جان و تن

شُعلۂ نمرود ہے روشن زمانے میں تو کیا

“شمع خود را می گدازد درمیانِ انجمن

نُورِ ما چوں آتشِ سنگ از نظر پنہاں خوش است”
</poem>
